\chapter{INTRODUCTION}
\label{ch:introduction}

% Purpose (from template): tell the reader WHY this study exists. Start broad,
% narrow to the research question. Define terms on first use. Reference prior work.
%
% Target length: ~600-900 words for this chapter (project is 5,000-10,000 total).
% See report/STRUCTURE.md for the full per-section budget.

Modern machine-learning systems represent text, images, and other media as
high-dimensional vectors called embeddings, arranged so that items with similar
meaning sit close together in the vector space. Searching a collection then
reduces to finding the stored vectors nearest to a query vector, a task known as
vector search. It underlies image search, recommendation systems, and the
retrieval step of retrieval-augmented language models. When the query and the
stored items are different kinds of media, such as retrieving images from a
written description, the task is cross-modal retrieval, and it becomes tractable
once a single embedding model places both modalities in one shared space.

Quantum computing offers two primitives that appear, at first sight, well matched
to this task. Grover's algorithm finds a marked item in an unstructured
collection using far fewer queries than a classical scan, and the swap test
estimates how similar two quantum states are. This study asks, through
measurement rather than assertion, whether those primitives are of practical use
for cross-modal vector search at the scale and on the hardware available today.
It is an investigation of where quantum search does and does not fit, not an
attempt to demonstrate a quantum advantage.

\section{Background of Study}
\label{sec:background}

% What is vector search, what is cross-modal retrieval, what is CLIP, and what
% does "quantum similarity" mean here. Two or three paragraphs. Cite the
% foundational papers (CLIP, Grover, Buhrman et al. swap test).

An embedding model maps an input to a fixed-length vector so that semantically
related inputs map to nearby vectors. For cross-modal retrieval we use CLIP, a
model trained on a large collection of image and caption pairs to place an image
and its description close together in a shared
space~\cite{radford2021clip}. Once both a text query and a set of images live in
that space, answering the query means ranking the image vectors by their
closeness to the query vector. The simplest way to do this is an exhaustive scan
that compares the query against every stored vector, which costs $O(N)$
comparisons for $N$ items. Production systems reduce this cost with specialised
libraries such as FAISS, which performs the exact scan
efficiently~\cite{johnson2019faiss}, or with approximate methods such as the
Hierarchical Navigable Small World (HNSW) graph, which reaches a near neighbour
in roughly $O(\log N)$ steps by navigating a layered proximity
graph~\cite{malkov2020hnsw}.

The quantum side of the comparison rests on two ideas. Grover's algorithm
searches an unstructured collection of $N$ items for a marked one using only
about $\sqrt{N}$ queries to an oracle, a quadratic improvement over the linear
classical scan~\cite{grover1996}. The swap test takes two quantum states and
estimates the squared overlap between them from a single measured
qubit~\cite{buhrman2001swaptest}; because a normalised embedding can be encoded
in the amplitudes of a quantum state, that overlap is a similarity score. Both
techniques therefore have a natural reading as vector-search operations, Grover
as the search itself and the swap test as the similarity comparison.

The difficulty is that these speed-ups assume the data is already available
inside the quantum computer. For Grover to search real vectors, an oracle would
have to recognise the closest item while the whole dataset is held in
superposition, which requires a quantum random-access memory that does not yet
exist at usable scale; loading the data one item at a time instead costs $O(N)$
and removes the very advantage the algorithm promises~\cite{aaronson2015readfine}.
This tension between an algorithm's asymptotic promise and the cost of getting
data into the machine is the central question the study examines.

\section{Problem Statement and Research Aims}
\label{sec:problem}

% Make the gap explicit: classical vector search is well-studied; the practical
% relevance of quantum vector search at NISQ scale is not. State what we set out
% to MEASURE - not what we set out to PROVE.

Classical vector search is mature, fast, and well understood. Whether quantum
methods add anything to it at the scale available today is far less clear, and
popular accounts often overstate the case. The gap this work addresses is
empirical: there are few honest, like-for-like comparisons that run classical and
quantum similarity search on the same data and report both retrieval accuracy and
the resources each method consumes. We built such a comparison. A single common
interface lets classical, quantum, and hybrid engines answer the same queries
over the same embeddings, so their rankings can be placed side by side and scored
identically.

We frame the work as a measurement rather than a proof of advantage. Present-day
devices are small and noisy, the regime usually called noisy intermediate-scale
quantum (NISQ), so we run the quantum engines mostly on a classical simulator and
validate the hardware path with a single small run on a real quantum processor.
The goal is to determine, on a controlled cross-modal dataset, how the quantum
and hybrid engines compare with strong classical baselines, and to state plainly
the conditions under which each approach is preferable.

\section{Research Questions and Objectives}
\label{sec:questions}

% List 3-4 measurable research questions. Each must map to something the
% benchmark harness already produces (MRR, oracle calls, circuit depth, etc.)
% so we never write a question we cannot answer with data we already have.

\begin{enumerate}[label=RQ\arabic*., leftmargin=*]
  \item \emph{Accuracy:} Does swap-test similarity reproduce the ranking of
        classical cosine similarity, and what is the shot budget required to do so?
  \item \emph{Scaling:} Does our Grover implementation reproduce the
        $\lfloor \pi\sqrt{N}/4 \rfloor$ oracle-call scaling in practice?
  \item \emph{Practicality:} Under what conditions \,--\, if any \,--\, would a
        quantum vector-search engine outperform a production classical baseline
        (HNSW) on the same dataset?
  \item \emph{Suitability:} Independent of hardware maturity, is cross-modal
        vector search a problem class that quantum computing is suited to, given
        that quantum advantage concentrates on compute-bound problems over small
        data rather than data-bound retrieval?
\end{enumerate}

To answer these questions we set four objectives, each corresponding to a
concrete artefact in the project repository:

\begin{enumerate}[leftmargin=*]
  \item Implement a suite of search engines spanning classical, quantum, and
        hybrid approaches behind a single common interface, so that one query can
        be answered by any of them and the results compared directly.
  \item Build a reproducible benchmark harness and a controlled cross-modal
        dataset, and collect a complete set of results across embedding
        dimensions and measurement shot counts.
  \item Provide a web application that displays the engines' rankings side by
        side for inspection.
  \item Produce an honest written analysis that answers the research questions
        from the measured data, including the cases where quantum methods do not
        help.
\end{enumerate}
